% Источник: lectures_08_10_sets.tex 118–123; lectures_07_10_random_variables.tex 10–35, 67–74, 93–112, 139–162, 273–281.
\question{Случайная величина и её распределение}

\begin{definition}[\topicmark{Борелевская сигма-алгебра}Определение]
    \textbf{Борелевской $\sigma$-алгеброй} на пространстве $A$
    называется наименьшая $\sigma$-алгебра, содержащая все
    открытые множества этого пространства.
    Обозначение: $\mathcal{B}(A)$.
\end{definition}

Пусть $(\Omega,\mathcal{F},\Pr)$ --- вероятностное пространство.
Идея: случайная величина --- числовая характеристика
случайного эксперимента.

\begin{definition}[\topicmark{Случайная величина}Определение]
    Функция $\xi\colon\Omega\to\R$ называется
    \textbf{случайной величиной}, если
    \[
        \xi^{-1}\bigl((-\infty;x]\bigr)\in\mathcal{F}
        \qquad\forall x\in\R.
    \]
    Эквивалентная запись:
    \[
        \{\omega\in\Omega\colon\xi(\omega)\leq x\}
        \in\mathcal{F}
        \qquad\forall x\in\R.
    \]
    Эквивалентное условие:
    \[
        \xi^{-1}(B)\in\mathcal{F}
        \qquad\forall B\in\mathcal{B}(\R).
    \]
    Здесь $\mathcal{B}(\R)$ --- наименьшая $\sigma$-алгебра,
    содержащая промежутки на $\R$, в частности
    $(y;x]$, $(y;x)$, $[y;x]$, $[y;x)$.
\end{definition}

\begin{definition}[\topicmark{Распределение}Определение]
    \textbf{Распределением случайной величины} $\xi$ называется
    вероятностная мера на $\mathcal{B}(\R)$, определённая равенством
    \[
        \Pr_\xi(B)=\Pr\{\omega\in\Omega\colon\xi(\omega)\in B\},
        \qquad B\in\mathcal{B}(\R).
    \]
\end{definition}

\begin{definition}[\topicmark{Дискретное распределение}Определение]
    Случайная величина $\xi$ имеет \textbf{дискретное распределение},
    если существует конечное или счётное множество $\{x_i\}_{i\in I}$,
    такое что
    \[
        \Pr\{\xi(\omega)\in\{x_i\}_{i\in I}\}=1.
    \]
    Распределение задаётся набором $\{(x_i;p_i)\}_{i\in I}$,
    где значения $x_i$ различны и
    \[
        p_i=\Pr\{\xi(\omega)=x_i\},\qquad
        p_i\geq0,\qquad \sum_{i\in I}p_i=1.
    \]
\end{definition}

Для дискретного распределения $\{(x_i;p_i)\}_i$:
\[
    \Pr_\xi(B)=\sum_{j\colon x_j\in B}p_j,
    \qquad B\in\mathcal{B}(\R).
\]

\begin{example}[\topicmark{Примеры}Три дискретных распределения]
\textbf{Бернулли.} В одном испытании $\xi\sim\operatorname{Bern}(p)$ равна $1$ с вероятностью $p$ и $0$ с вероятностью $1-p$.

\smallskip
\textbf{Биномиальное.} Число успехов в $n$ испытаниях Бернулли имеет распределение $\operatorname{Bin}(n;p)$:
\[
 \Pr\{\xi=k\}=C_n^k p^k(1-p)^{n-k},\qquad k=0,\ldots,n.
\]
\textbf{Геометрическое.} Номер первого успеха в бесконечной схеме Бернулли имеет распределение $\operatorname{Geom}(p)$:
\[
 \Pr\{\xi=k\}=(1-p)^{k-1}p,\qquad k=1,2,\ldots
\]
\end{example}

% Источник: lectures_07_10_random_variables.tex 192–269.
\question{Пуассоновское распределение. Теорема Пуассона}

\begin{definition}[\topicmark{Распределение Пуассона}Определение]
    \[
        \operatorname{Poiss}(\lambda),\qquad\lambda>0.
    \]
    \[
        x_k=k\in\{0,1,2,\ldots\},\qquad
        p_k=\frac{\lambda^k}{k!}e^{-\lambda}.
    \]
\end{definition}

\begin{theorem}[\topicmark{Теорема Пуассона}Теорема Пуассона]
    Рассмотрим серию схем испытаний Бернулли:
    \[
        \xi_n\sim\operatorname{Bin}(n;p_n),\qquad
        np_n\xrightarrow[n\to\infty]{}\lambda,\qquad
        \eta\sim\operatorname{Poiss}(\lambda).
    \]
    Тогда для каждого фиксированного $k=0,1,2,\ldots$
    \[
        \Pr\{\xi_n=k\}\xrightarrow[n\to\infty]{}\Pr\{\eta=k\}.
    \]
\end{theorem}

\begin{Proof}
    \[
        \begin{aligned}
            \Pr\{\xi_n=k\}
            &=C_n^k p_n^k(1-p_n)^{n-k}\\
            &=\frac{n!}{(n-k)!\,k!}\,p_n^k(1-p_n)^{n-k}\\
            &=\frac1{k!}\,
              \underbrace{\frac{n(n-1)\cdots(n-k+1)}{n^k}}
              _{\longrightarrow\,1}\,
              \underbrace{(np_n)^k}_{\longrightarrow\,\lambda^k}\,
              \frac{(1-p_n)^n}
              {\underbrace{(1-p_n)^k}_{\longrightarrow\,1}}.
        \end{aligned}
    \]
    При этом
    \[
        (1-p_n)^n
        =\left(\underbrace{(1-p_n)^{-1/p_n}}
        _{\longrightarrow\,e}\right)^{-np_n}
        \longrightarrow e^{-\lambda}.
    \]
    Следовательно,
    \[
        \Pr\{\xi_n=k\}\longrightarrow
        \frac1{k!}\cdot1\cdot\lambda^k e^{-\lambda}
        =\frac{\lambda^k}{k!}e^{-\lambda}.
    \]
\end{Proof}

\begin{theorem}[\topicmark{Обобщение}Обобщение]
    Пусть $\xi\sim\operatorname{Bin}(n;p)$,
    $\eta\sim\operatorname{Poiss}(\lambda)$, $np=\lambda$.
    Тогда
    \[
        \bigl|\Pr\{\xi\in B\}-\Pr\{\eta\in B\}\bigr|
        \leq np^2
        \qquad\forall B\in\mathcal{B}(\R).
    \]
\end{theorem}

% Источник: lectures_07_10_random_variables.tex 302–309, 343–386.
\question{Функция распределения и её свойства}

\begin{definition}[\topicmark{Определение}Определение]
    \textbf{Функцией распределения случайной величины} $\xi$
    называется функция
    \[
        F_\xi(x)=\Pr\{\xi(\omega)\leq x\}
        =\Pr_\xi\bigl((-\infty;x]\bigr),\qquad x\in\R.
    \]
\end{definition}

\begin{theorem}[\topicmark{Свойства}Теорема]
    \begin{enumerate}[label=\arabic*)]
        \item $\lim_{x\to+\infty}F_\xi(x)=1$,
        $\lim_{x\to-\infty}F_\xi(x)=0$.
        \item $F_\xi$ не убывает.
        \item $F_\xi$ непрерывна справа.
    \end{enumerate}
\end{theorem}

\begin{Proof}[Монотонность и пределы\editorialmark]
При $x\leq y$ вложение $\{\xi\leq x\}\subseteq\{\xi\leq y\}$ даёт $F_\xi(x)\leq F_\xi(y)$.

\smallskip
Поскольку $\xi$ принимает конечные вещественные значения, при $n\to\infty$
\[
  \{\xi\leq n\}\uparrow\Omega,\qquad
  \{\xi\leq -n\}\downarrow\varnothing.
\]
По непрерывности вероятности снизу и сверху $F_\xi(n)\to 1$ и $F_\xi(-n)\to 0$.
Монотонность даёт те же пределы при $x\to+\infty$ и $x\to-\infty$ соответственно.
\end{Proof}

\begin{Proof}[\topicmark{Непрерывность справа}Непрерывность справа]
    По Гейне рассмотрим $x_n\downarrow x_0$. Положим
    \[
        A_n=\xi^{-1}\bigl((-\infty;x_n]\bigr),
        \qquad A=\xi^{-1}\bigl((-\infty;x_0]\bigr).
    \]
    Тогда
    \[
        F_\xi(x_n)=\Pr\{\xi(\omega)\leq x_n\}=\Pr(A_n),
        \qquad F_\xi(x_0)=\Pr(A).
    \]
    Поскольку $A_n\downarrow A$, по непрерывности сверху
    \[
        \Pr(A_n)\longrightarrow\Pr(A),
        \qquad F_\xi(x_n)\longrightarrow F_\xi(x_0).
    \]
    Следовательно, $\lim_{x\to x_0+}F_\xi(x)=F_\xi(x_0)$.
\end{Proof}

\begin{remark}
    Для неубывающей, ограниченной и непрерывной справа функции $F$
    на полуинтервалах задаётся
    \[
        m\bigl((a;b]\bigr)=F(b)-F(a).
    \]
\end{remark}

\begin{theorem}[\topicmark{Следствие из теоремы Каратеодори}Следствие из теоремы Каратеодори]
    Существует взаимно однозначное соответствие:
    \[
        \text{вероятностные меры на }\mathcal{B}(\R)
        \quad\longleftrightarrow\quad
        \text{функции распределения}.
    \]
\end{theorem}

% ТВ 3.pdf, с. 5; Тервер(Лекция)68.pdf, с. 12–13.
На полукольце мера задана равенством $\Pr_\xi((a;b])=F_\xi(b)-F_\xi(a)$.
По теореме Каратеодори её продолжение на $\mathcal B(\R)$ единственно.

% Источник: lectures_07_10_random_variables.tex 390–434.
\question{Непрерывное и абсолютно непрерывное распределения}

\begin{definition}[\topicmark{Непрерывное распределение}Непрерывное распределение\editorialmark]
Распределение $\xi$ называется \textbf{непрерывным}, если его функция
распределения $F_\xi$ непрерывна в каждой точке $\R$.

\smallskip
Равносильное условие: $\Pr\{\xi=x\}=0$ для каждого $x\in\R$.
Действительно, $F_\xi(x-)=\Pr\{\xi<x\}$, поэтому
\[
 F_\xi(x)-F_\xi(x-)=\Pr\{\xi=x\}.
\]
Справа функция распределения непрерывна всегда.
\end{definition}

\begin{definition}[\topicmark{Плотность}Определение]
    Распределение называется \textbf{абсолютно непрерывным},
    если существует неотрицательная измеримая функция
    $p_\xi\colon\R\to[0;+\infty)$, называемая
    \textbf{плотностью распределения}, такая что
    \[
        p_\xi(x)\geq0,\qquad
        F_\xi(x)=\int_{-\infty}^{x}p_\xi(t)\,\diff t.
    \]
    При этом
    \[
        \int_{-\infty}^{+\infty}p_\xi(t)\,\diff t=1.
    \]
    Эквивалентно, \textbf{плотность распределения} $p_\xi$ задаёт
    вероятности попадания в борелевские множества:
    \[
        \Pr\{\omega\in\Omega:\xi(\omega)\in B\}
        =\int_B p_\xi(x)\,\diff x,
        \qquad B\in\mathcal{B}(\R).
    \]
\end{definition}

\begin{definition}[\topicmark{Абсолютная непрерывность}Абсолютная непрерывность функции распределения]
    Функция распределения $F$ \textbf{абсолютно непрерывна}, если
    для каждого $\eps>0$ существует $\delta>0$ такое, что для любого
    конечного семейства попарно непересекающихся интервалов $(a_i;b_i)$
    \[
        \sum_{i=1}^{m}(b_i-a_i)<\delta
        \quad\Longrightarrow\quad
        \sum_{i=1}^{m}\bigl(F(b_i)-F(a_i)\bigr)<\eps.
    \]
    Для функции распределения это условие равносильно существованию
    плотности. Непрерывности $F$ самой по себе недостаточно.
\end{definition}

\begin{remark}[\topicmark{Плотность и производная}Плотность и производная функции распределения]
    Плотность определяется однозначно с точностью до множества
    меры Лебега нуль: изменение её значений на таком множестве
    не меняет распределение.

    \smallskip
    Для абсолютно непрерывной функции распределения
    \[
        F_\xi'(x)=p_\xi(x)
        \qquad\text{почти всюду относительно меры Лебега}.
    \]
\end{remark}

% Источник: lecture_11_distributions.tex 8–90.
\question{Равномерное, экспоненциальное и нормальное распределения}

\begin{definition}[\topicmark{Равномерное}Равномерное распределение]
    Случайная величина $\xi$ имеет \textbf{равномерное распределение}
    на отрезке $[a;b]$, где $a<b$, если её плотность имеет вид
    \[
        p_\xi(t)=\frac{1}{b-a}\,I_{[a;b]}(t)
        =\begin{cases}
            \dfrac{1}{b-a},&a\leq t\leq b,\\
            0,&t\notin[a;b].
        \end{cases}
    \]
    Обозначение: $\xi\sim\Uniform[a;b]$.
    Здесь $I_A(t)=1$ при $t\in A$ и $I_A(t)=0$ при $t\notin A$.
\end{definition}

\begin{definition}[\topicmark{Экспоненциальное}Экспоненциальное распределение]
    Случайная величина $\xi$ имеет \textbf{экспоненциальное распределение}
    с параметром $\lambda>0$, если её плотность имеет вид
    \[
        p_\xi(t)=\lambda e^{-\lambda t}I_{[0;+\infty)}(t)
        =\begin{cases}
            \lambda e^{-\lambda t},&t\geq0,\\
            0,&t<0.
        \end{cases}
    \]
    Обозначение: $\xi\sim\operatorname{Exp}(\lambda)$.
\end{definition}

\begin{definition}[\topicmark{Нормальное}Нормальное распределение]
    Случайная величина $\xi$ имеет \textbf{нормальное распределение}
    с параметрами $a\in\R$ и $\sigma^2>0$, если её плотность равна
    \[
        p_\xi(t)=\frac{1}{\sqrt{2\pi}\,\sigma}
        \exp\left(-\frac{(t-a)^2}{2\sigma^2}\right),
        \qquad t\in\R,\quad\sigma>0.
    \]
    Обозначение: $\xi\sim N(a;\sigma^2)$.
    Распределение $N(0;1)$ называется \textbf{стандартным нормальным};
    его плотность равна
    \[
        p_\xi(t)=\frac{1}{\sqrt{2\pi}}e^{-t^2/2}.
    \]
\end{definition}

% Источник: lecture_11_distributions.tex 93–152.
\question{Интеграл Гаусса. Нормировка нормальной плотности}

\begin{lemma}[\topicmark{Интеграл Гаусса}Интеграл Гаусса]
    \[
        \int_{-\infty}^{+\infty}e^{-x^2}\,\diff x=\sqrt\pi.
    \]
\end{lemma}

\begin{Proof}[\topicmark{Доказательство}Доказательство]
    Обозначим интеграл через $C$. Он положителен и конечен:
    при $|x|\geq1$ имеем $e^{-x^2}\leq e^{-|x|}$.
    Тогда
    \[
        \begin{aligned}
            C^2
            &=\left(\int_{-\infty}^{+\infty}e^{-x^2}\,\diff x\right)
              \left(\int_{-\infty}^{+\infty}e^{-y^2}\,\diff y\right)\\
            &=\int_{-\infty}^{+\infty}\int_{-\infty}^{+\infty}
              e^{-(x^2+y^2)}\,\diff x\,\diff y
              =\iint_{\R^2}e^{-(x^2+y^2)}\,\diff x\,\diff y.
        \end{aligned}
    \]
    Перейдём к полярным координатам:
    \[
        x=r\cos\varphi,\qquad y=r\sin\varphi,
        \qquad r\geq0,\quad0\leq\varphi<2\pi.
    \]
    Якобиан замены равен
    \[
        \left|\frac{\partial(x,y)}{\partial(r,\varphi)}\right|
        =\left|\begin{vmatrix}
            \cos\varphi&-r\sin\varphi\\
            \sin\varphi&r\cos\varphi
        \end{vmatrix}\right|=r.
    \]
    Поэтому
    \[
        \begin{aligned}
            C^2&=\int_0^{+\infty}\int_0^{2\pi}
                  e^{-r^2}r\,\diff\varphi\,\diff r\\
               &=2\pi\int_0^{+\infty}r e^{-r^2}\,\diff r
                 =\pi\int_0^{+\infty}e^{-u}\,\diff u=\pi,
                 \qquad u=r^2.
        \end{aligned}
    \]
    Так как $C>0$, получаем $C=\sqrt\pi$.
\end{Proof}

\begin{Proof}[\topicmark{Нормировка плотности}Нормировка нормальной плотности]
    При замене $y=(t-a)/(\sqrt2\,\sigma)$ имеем
    $\diff t=\sqrt2\,\sigma\,\diff y$. Следовательно,
    \[
        \begin{aligned}
            \int_{-\infty}^{+\infty}p_\xi(t)\,\diff t
            &=\frac{1}{\sqrt{2\pi}\,\sigma}
              \int_{-\infty}^{+\infty}
              e^{-(t-a)^2/(2\sigma^2)}\,\diff t\\
            &=\frac1{\sqrt\pi}\int_{-\infty}^{+\infty}
              e^{-y^2}\,\diff y=1.
        \end{aligned}
    \]
\end{Proof}

% Источник: lecture_17_distribution_decomposition.tex 8–28.
\question{Сингулярные распределения. Разложение функции распределения}

\begin{definition}[\topicmark{Сингулярное распределение}Сингулярное непрерывное распределение]
    Функция распределения $F$ называется \textbf{сингулярной непрерывной},
    если она непрерывна и
    \[
        F'(x)=0
        \quad\text{почти всюду относительно меры Лебега}.
    \]
\end{definition}

\begin{theorem}[\topicmark{Теорема о разложении}Разложение функции распределения]
    Любую функцию распределения можно представить в виде
    \[
        F(x)=\alpha_1F_1(x)+\alpha_2F_2(x)+\alpha_3F_3(x),
    \]
    где $F_1$ --- дискретная, $F_2$ --- абсолютно непрерывная,
    $F_3$ --- сингулярная непрерывная функция распределения, причём
    \[
        \alpha_1,\alpha_2,\alpha_3\geq0,
        \qquad \alpha_1+\alpha_2+\alpha_3=1.
    \]
\end{theorem}

